\section{Signed rows, zero extensions, and label reconstruction}
\label{sec:signed-reconstruction}

The purpose of the reconstruction is to identify one integer row to which
the unmarked density theorem applies. It must identify the character on
all arguments, including its zeros, and must also explain which arithmetic
labels can be recovered from that row. Merely identifying a primitive
character on coprime arguments would lose the natural masks.

\subsection{The two signs at a common prime}
For supported primary elements define
$c_n(v)=\mathfrak s_v((n))$, the source's zero-extended sextic ideal row.
Let $q(\pi,n)$ be the fixed reciprocity phase, so
\[
 c_\pi(n)=q(\pi,n)c_n(\pi).
\]
The phase belongs to the fixed ray dictionary. If $\pi\mid n$, both
sides vanish; the identity is not restricted to the unit case.
For a common prime with active net exponent $1\le t\le5$, it follows that
\begin{align*}
 c_n(\pi)c_\pi(n)^t
   &=q(\pi,n)^t c_n(\pi)^{1+t},\\
 c_n(\pi)\overline{c_\pi(n)^t}
   &=q(\pi,n)^{6-t}c_n(\pi)^{7-t}.
\end{align*}
On units this is ordinary multiplication of sixth roots of unity;
off units each displayed product is zero. Thus the left coefficient has
residue $1+t$ and the right coefficient residue $1-t$, modulo six.
These are the signs in the actual formal coefficients
\path{leftCoefficient} and \path{rightCoefficient}.
A manuscript convention assigning the conjugate coefficient to the
C-side must exchange the side names consistently.

If $t=5$ on the left, or $t=1$ on the right, the resulting exponent is six.
The source's exact identity is
\[
 c_n(\pi)^6=\mathbf1_{(n,\pi)=1}.
\]
Consequently reducing the exponent to zero moves this factor into the
puncture; it does not replace it by one. Multiplying the primewise
identities proves the corresponding complete first-row character
identity with every canceled prime still punctured.

For one side let $F=\prod_{\pi\mid r}\pi^{e(\pi)}$, choosing the
source's primary prime generators, and let $L=\log_Z NF$. The ideal
used by the formal source has generator
$\operatorname{idealGenerator}(\mathfrak r)$, which can differ from the
product of chosen generators by a unit. Retain that unit in the effective
row, or absorb its character into the fixed six-unit dictionary. No
unproved equality of these two generator choices is needed.

\subsection{The complete second frequency}
For the second common-support ideals $D_2,E_2$, let $G$ be their common
gcd generator. The actual correlation vanishes unless $G\mid j$.
After writing $j=Gu$, the nonunit part of the equal-multiplicity
partition forces $V\mid u$, where $V$ is its radical generator.
Thus the genuine nonzero frequency is
\[
 j=GVh',\qquad h'\ne0.
\]
The source's partitioned correlation, divided by its stated normalizer,
has modulus at most one, including its extension by zero off the
partition. Multiplicativity now combines the first row with this frequency:
\[
 w=u_0\,eFGVh',
\]
where $u_0$ is the retained unit. Every surviving whole-product character
is a fixed ray factor times $c_n(w)$, with the common puncture and all
canceled-prime holes. Relative slot characters remain their original
fixed factors. Whole-character conjugation or the reflected row $-w$
is handled by whole-character height reflection and the fixed unit
list. The dictionary is a finite union of parent translates of the
allowed $Q_*$-character group. The buffered maximum is formed within
the actual parent translate; another translate is not added to that
maximum. Conjugating only its fixed factor would change this family
and is not an allowed replacement.

The exact source ingredients are
\path{SecondCanonicalFrequency.idealCorrelation_zero_outside_common},
\path{SecondCanonicalNonunit.nonunitFrequencyGenerator_dvd}, and
\path{SecondCanonicalFrequency.actual_ideal_partition_scalar}.
The new signed-row lemmas check the primewise identities and their
multiplicative combination against these source conventions.

With the actual common main ball of logarithmic radius
$K_0+(d-c)_++6a+\varpi$, the second frequency envelope gives
\[
 \log_Z Nw\le M-R-(c-d)_++L-6a-\varpi+\eta.
\]
Discarding the favorable $-\varpi$ gives \eqref{eq:integerrow}.
A nonempty integral row window always has $Nw\ge1$; negative nominal
logarithmic ranges are either empty or are filled by the reserved
bounded-window allowance. Density is not applied as a count smaller
than one on a nonempty window.

\subsection{Which labels are divisors of the row?}
The first residual projector label $s$ divides both second columns.
After the complete second common-support extraction it therefore divides
$G$, and hence divides $w$. The entire second common radical also divides
$G$. The source proves these assertions as
\path{SecondSourceMask.source_pair_divisor_common} and
\path{divisor_dvd_common}.

The second residual projector label $t$ is different. Its individual
terms belong to the artificial separated extension of the genuine
correlation, and $t$ need not divide $G$ or $w$. Freeze it while estimating
the row norm. Its signed priority allocation selects plain divisors and
freezes prime slots. If a selected plain is $n=dv$, then
\[
 \psi_w(dv)=\psi_w(d)\psi_w(v).
\]
The factor $\psi_w(d)$ is retained outside, even if it is zero. The
quotient $v$ may still be divisible by a prime of $d$; no new coprimality
condition is introduced. This explains why the primitive row and its
buffered bin survive allocation. Its divisor-dyad count is paid together
with the two exact normalization scalars in B3.67--B3.69.

Among the first common primes, those active with nonzero row residue
occur in $F$, and the selected mask-divisor primes occur in $e$.
The remaining first radical has logarithmic width at most
\[
 p-R-E+R_{\rm zero}.
\]
This accounts separately for inactive net-zero primes and for active
primes whose row phase cancels. They cannot be called divisors of $w$.
Once their radical is fixed, all remaining first labels are supported
on it together with $\operatorname{rad}(w)$.

For completeness, a uniform support-counting bound follows from Rankin's
inequality. If $NQ\le Z^B$ and $NI\le Z^T$, then for any $s>0$,
\[
 \#\{I:\operatorname{rad}I\mid Q,\ NI\le Z^T\}
 \le Z^{sT}\prod_{\mathfrak p\mid Q}(1-(N\mathfrak p)^{-s})^{-1}.
\]
Choose $sT$ small. Each sufficiently large prime factor in the product
is at most $(N\mathfrak p)^b$ for any prescribed $b>0$; the finitely
many remaining primes contribute a constant. Choosing $bB$ small proves
a bound $O(Z^\epsilon)$ uniformly in the support ideal. The field has
only $O(X)$ ideals of norm at most $X$, so counting the missing radical
costs its displayed norm exponent. Fixed finitely many slot assignments,
bounded-range powers and priority allocations have the same arbitrarily
small power allowance after the slot system is fixed.

This proves the general label reconstruction used in the high norm.
It retains the first labels until the joint row count, and it does not
charge the second projector count twice. Nonprincipality and primitive
conductor bounds are supplied by the following fixed-dictionary section;
they are not consequences of the size bound on $w$ alone.
